\documentclass[a4paper,10pt]{article}
\usepackage[utf8]{inputenc}
\usepackage{adjustbox}
\usepackage{float}
\usepackage{indentfirst}
\usepackage{hyperref}
\usepackage{pifont}
\newcommand{\cmark}{\ding{51}}%
\newcommand{\xmark}{\ding{55}}%

%opening
\title{
	$SiER^2NA$ \\
	\small Simple Extensible Register-Register Neumann Architecture \\
	\large Instruction Set Architecture \\
	\large 1.0
}
\date{\today}

\begin{document}

\maketitle
\tableofcontents

\section{Preface}
\section{Glossary}
\begin{itemize}
	\item Word - 32-bit value
	\item Half-word - 16-bit value
	\item Byte - 8-bit value
\end{itemize}

\section{Registers}

The SyrENA architecture, similarily to the TMS9900, does not store general-purpose registers on the CPU, but rather a \emph{frame pointer} that points to the address in register RAM of the current \emph{register frame} for the executing context.

A \emph{register frame} refers to a (virtual) structure in RAM of 128 bytes (1024 bits) in length that stores 32 4-byte values total:

\begin{table}[H]
	\caption{Register frame (1023-768)}
	\centering
	\resizebox{\linewidth}{!}{
		\begin{tabular}{|c|c|c|c|c|c|c|c|}
\hline
1023..992 & 991..960 & 959..928 & 927..896 & 895..864 & 863..832 & 831..800 & 799..768 \\
\hline
rfo1 & rfo0 & rfi1 & rfi0 & \multicolumn{4}{|c|}{reserved} \\
\hline
		\end{tabular}
	}
\end{table}

\begin{table}[H]
	\caption{Register frame (767-512)}
	\centering
	\resizebox{\linewidth}{!}{
		\begin{tabular}{|c|c|c|c|c|c|c|c|}
\hline
767..736 & 735..704 & 703..672 & 671..640 & 639..608 & 607..576 & 575..544 & 543..512 \\
\hline
\multicolumn{4}{|c|}{reserved} & rlfp & rsp & rtp & rip \\
\hline
		\end{tabular}
	}
\end{table}

\begin{table}[H]
	\caption{Register frame (511-256)}
	\centering
	\resizebox{\linewidth}{!}{
		\begin{tabular}{|c|c|c|c|c|c|c|c|}
\hline
511..480 & 479..448 & 447..416 & 415..384 & 383..352 & 351..320 & 319.288 & 287..256 \\
\hline
r15 & r14 & r13 & r12 & r11 & r10 & r9 & r8 \\
\hline
		\end{tabular}
	}
\end{table}

\begin{table}[H]
	\caption{Register frame (255-0)}
	\centering
	\resizebox{\linewidth}{!}{
		\begin{tabular}{|c|c|c|c|c|c|c|c|}
\hline
255..224 & 223..192 & 191..160 & 159..128 & 127..96 & 95..64 & 63..32 & 31..0 \\
\hline
r7 & r6 & r5 & r4 & r3 & r2 & r1 & r0 \\
\hline
		\end{tabular}
	}
\end{table}

Registers are categorised into two categories:
\begin{itemize}
	\item \textbf{Global} registers are stored directly on the CPU and their values persist between frame changes
	\item \textbf{Local} registers are stored in the \emph{register frame} and their address and/or value depends on the current \emph{frame pointer}
\end{itemize}

\subsection{Pointer registers}
\subsubsection{rfp - frame pointer}
{\small \textit{Availability: Global}}

Stores the address to the current \emph{register frame}

\subsubsection{rrap - interrupt routine array pointer}
{\small \textit{Availability: Global}}

Stores the address to the current \hyperref[ir:ARRAY]{Interrupt Routine Array}. Set by the system to indicate which IRA the CPU should use for interrupt handling

\subsubsection{rlfp - last frame pointer}
{\small \textit{Availability: Local}}

Stores the address to the previous \emph{register frame}. Used when returning from interrupt routines

\subsubsection{rip - instruction pointer}
{\small \textit{Availability: Local}}

Stores the address of the currently executing instruction

\subsubsection{rtp - procedure tail pointer}
{\small \textit{Availability: Local}}

Stores the address of the instruction that called the current procedure

\subsubsection{rsp - stack pointer}
{\small \textit{Availability: Local}}

Stores the address of the current stack frame

\subsection{Restricted purpose}
\subsubsection{rz - zero register}
{\small \textit{Availability: Global}}

Always 0. Writes are ignored

\subsection{General purpose}
\subsubsection{r0}
{\small \textit{Availability: Local}}

\subsubsection{r1}
{\small \textit{Availability: Local}}

\subsubsection{r2}
{\small \textit{Availability: Local}}

\subsubsection{r3}
{\small \textit{Availability: Local}}

\subsubsection{r4}
{\small \textit{Availability: Local}}

\subsubsection{r5}
{\small \textit{Availability: Local}}

\subsubsection{r6}
{\small \textit{Availability: Local}}

\subsubsection{r7}
{\small \textit{Availability: Local}}

\subsubsection{r8}
{\small \textit{Availability: Local}}

\subsubsection{r9}
{\small \textit{Availability: Local}}

\subsubsection{r10}
{\small \textit{Availability: Local}}

\subsubsection{r11}
{\small \textit{Availability: Local}}

\subsubsection{r12}
{\small \textit{Availability: Local}}

\subsubsection{r13}
{\small \textit{Availability: Local}}

\subsubsection{r14}
{\small \textit{Availability: Local}}

\subsubsection{r15}
{\small \textit{Availability: Local}}

\subsection{Special registers}
\subsubsection{rfi0}
{\small \textit{Availability: Local}}

Used for input flags for certain operations

\subsubsection{rfi1}
{\small \textit{Availability: Local}}

Used for input flags for certain operations

\subsubsection{rfo0}
{\small \textit{Availability: Local}}

Used for output flags of certain operations

\subsubsection{rfo1}
{\small \textit{Availability: Local}}

Used for output flags of certain operations

\section{Interrupts}
Interrupts are handled by the CPU after every execute cycle of a single instruction. The CPU checks if any interrupts were raised after executing the instruction and, if so, checks and picks the correct \hyperref[ir:HANDLER]{Interrupt Routine Handler} from the current \hyperref[ir:ARRAY]{Interrupt Routine Array} to handle the raised interrupt.

If the \verb|enable| bit of the IRA entry is not set, the interrupt is not handled in software, but the exact behaviour depends on the type of the raised interrupt. \\

There are three types of interrupts:

\begin{table}[H]
	\centering
	%\resizebox{\linewidth}{!}
	{
		\begin{tabular}{|c|c|c|c|c|}
\hline
Type & SW & HW & Recoverable & Pauses execution \\
\hline
Signal & \cmark & \cmark & \cmark & \xmark \\
\hline
Trap & \cmark & \cmark & \cmark & \cmark \\
\hline
Fault & \xmark & \cmark & \xmark & \cmark \\
\hline
		\end{tabular}
	}
\end{table}

\begin{itemize}
	\item \emph{Signal} interrupts do not pause execution and do not have to be explicitly handled in any way. This type of interrupt is recommended for general hardware-software communication.
	\item Raised \emph{trap} interrupts are most commonly recoverable errors and will pause execution of the frame that the trap occurred in if no valid IRH was found to handle the raised interrupt.
	\item Raised \emph{fault} interrupts are non-recoverable and reset the CPU upon existing from the IRH or if no valid handler was found. These cannot be raised manually from software.
\end{itemize}

Upon startup the \emph{power-on program} must set the \verb|rrap| register to a new empty, zero-filled IRA. Upon taking control, the operating system should allocate its own IRA, filling it with whatever entries are necessary to handle wanted interrupts, and set the value of \verb|rrap| to the address of said IRA.

\subsection{Interrupt List}
\begin{table}[H]
	\centering
	%\resizebox{\linewidth}{!}
	{
		\begin{tabular}{|c|c|c|c|c|}
\hline
ID & Mnemonic & Name & Type & Notes \\
\hline
0x00 & \#UDO & Undefined Opcode & Trap & \\
\hline
0x01 & \#DBG & Debug & Trap & \\
\hline
0x02 & \#BRK & Breakpoint & Trap & \\
\hline
0x03 & \#IOP & Invalid Operation & Trap & See \ref{ir:IOP} \\
\hline
0x04 & \#DBZ & Division By Zero & Trap & \\
\hline
0x05 & & & & \\
.. & & & & reserved \\
0x1F & & & & \\
\hline
0x20 & & & & \\
.. & & & & user-defined \\
0xFF & & & & \\
\hline
		\end{tabular}
	}
\end{table}

\subsubsection{\#IOP Invalid Operation} \label {ir:IOP}
The \emph{Invalid Operation} interrupt sets the \verb|rfo0| register accordingly based on the type of the operation that caused the interrupt:

\begin{table}[H]
	\centering
	{
		\begin{tabular}{|c|c|c|}
\hline
rfo0 & Operation \\
\hline
% TODO
		\end{tabular}
	}
\end{table}

\subsection{Interrupt Routine Array} \label{ir:ARRAY}
\begin{table}[H]
	\centering
	{
		\begin{tabular}{|c|c|}
\hline
Offset & Value \\
\hline
0x00 & Entry for INT 0x00 \\
\hline
0x0C & Entry for INT 0x01 \\
\hline
0x18 & Entry for INT 0x02 \\
\hline
... & ... \\
\hline
0xBF4 & Entry for INT 0xFF \\
\hline
\multicolumn{2}{|c|}{0xC00} \\
\hline
		\end{tabular}
	}
\end{table}

\subsubsection{Interrupt Routine Array Entry}
\begin{table}[H]
	\centering
	{
		\begin{tabular}{|c|c|c|c|}
\hline
95..64 & 63..32 & 31..1 & 0 \\
\hline
\&rp & \&rfp & -- & en \\
\hline
Interrupt Routine Handler & Register Frame & Reserved & Enabled \\
\hline
		\end{tabular}
	}
\end{table}

\subsection{Interrupt Routine Handler} \label{ir:HANDLER}
Interrupt Routine Handlers (IRH) are user-defined procedures that respond to a specific interrupt raised by the CPU. IRs have their own \emph{register frame} and can access the previous \emph{register frame} via its address located in the \verb|rlfp| register.

IRHs can exit at any time and resume execution by calling \verb|reti|. For non-recoverable interrupts this will result in the CPU being reset.

\subsection{User-defined interrupts}
% TODO

\section{Memory management}
\subsection{Paging}
The SiERNA architecture utilises \emph{paging} as its technique for memory protection and virtual address space mapping. A single page spans \textbf{4 KiB} (0x1000) of physical memory.

The page table is a 3-level structure; the full 32-bit physical address space (0xFFFF FFFF) is encapsulated by a level-1 page directory which contains 256 entries, each pointing to a level-2 page directory (or 0). A level-2 page directory contains 256 entries, each pointing to a page table (or 0). Every page table entry (page descriptor) then describes the page as well as its physical address.

% split into 256 (0x100) level-1 page directories, each of switch addresses a level-2 page directory. A level-2 page directory then stores 256 entries, each addressing a page table. Finally, each page table addresses 16 individual pages (page descriptors).

\subsubsection{Page Directory}
\begin{table}[H]
	\centering
	{
		\begin{tabular}{|c|c|}
\hline
Offset & Value \\
\hline
0x00 & Entry 0 \\
\hline
0x04 & Entry 1 \\
\hline
0x08 & Entry 2 \\
\hline
... & ... \\
\hline
0x3FC & Entry 255 \\
\hline
\multicolumn{2}{|c|}{0x400} \\
\hline
		\end{tabular}
	}
\end{table}

\subsubsection{Page Directory Entry}
\begin{table}[H]
	\centering
	{
		\begin{tabular}{|c|}
\hline
31..0 \\
\hline
\&adr \\
\hline
Physical page directory/table address \\
\hline
		\end{tabular}
	}
\end{table}

\subsubsection{Page Table}
\begin{table}[H]
	\centering
	{
		\begin{tabular}{|c|c|}
\hline
Offset & Value \\
\hline
0x00 & Descriptor 0 \\
\hline
0x08 & Descriptor 1 \\
\hline
0x10 & Descriptor 2 \\
\hline
... & ... \\
\hline
0x78 & Descriptor 15 \\
\hline
\multicolumn{2}{|c|}{0x80} \\
\hline
		\end{tabular}
	}
\end{table}

\subsubsection{Page Descriptor}
\begin{table}[H]
	\centering
	{
		\begin{tabular}{|c|c|}
\hline
31..12 & 11..0 \\
\hline
\&adr & -- \\
\hline
Physical page address & Reserved \\
\hline
		\end{tabular}
	}
\end{table}

% \subsection{Partitioning}
% The SiERNA architecture utilises so-called \emph{memory partitioning} for protection and virtual address spaces. It can be thought of as an equivalent to segmentation on the x86.
%
% \subsubsection{Partition Descriptor Node}
% \begin{table}[H]
% 	\centering
% 	{
% 		\begin{tabular}{|c|c|}
% \hline
% 63..32 & 31..0 \\
% \hline
% \&nx & \&vl \\
% \hline
% Next & Value \\
% \hline
% 		\end{tabular}
% 	}
% \end{table}
%
% Value points to the PD that this node describes.
%
% Next points to the next value. Can be 0.
%
% \subsubsection{Partition Descriptor}
% \begin{table}[H]
% 	\centering
% 	{
% 		\begin{tabular}{|c|c|c|}
% \hline
% 95..64 & 63..32 & 31..0 \\
% \hline
% -- & \&li & \&bs \\
% \hline
% Reserved & Limit & Base \\
% \hline
% 		\end{tabular}
% 	}
%\end{table}

\subsection{Virtual memory}

\section{Structures}
\subsection{CPUID} \label{st:CPUID}

\begin{table}[H]
	\centering
	\resizebox{\linewidth}{!}{
		\begin{tabular}{|l|l|l|l|l|l|l|l|l|l|l|l|l|l|l|l|l|l|l|l|l|l|l|l|l|l|l|l|l|l|l|l|}
\hline
31 & 30 & 29 & 28 & 27 & 26 & 25 & 24 & 23 & 22 & 21 & 20 & 19 & 18 & 17 & 16 & 15 & 14 & 13 & 12 & 11 & 10 & 9 & 8 & 7 & 6 & 5 & 4 & 3 & 2 & 1 & 0 \\
\hline
\multicolumn{16}{|c|}{reserved} & \multicolumn{16}{|c|}{available instruction sets} \\
\hline
		\end{tabular}
	}
\end{table}

\end{document}
